\chapter{Hereditary Harrop formulas: implication and universal goals}\label{ch:hh}

% local notation: a focused (backchaining) sequent  Σ; P --D--> A
\newcommand{\hhfoc}[4]{#1;#2\xrightarrow{\,#3\,}#4}

In \cref{ch:horn} we met the logic of first-order Horn clauses, \fohc, the logic behind Prolog.
Its search semantics was given by a handful of reduction rules, one per logical symbol, and two
of them never fired: AUGMENT, for goals of the form $D\imp G$, and GENERIC, for goals of the form
$\forall x\,G$. The reason is purely syntactic: in \fohc{} neither goals nor clause bodies may
contain implications or universal quantifiers, so the signature $\Sig$ and the program $\Prog$ stay
fixed for the whole computation. Miller and Nadathur \cite{MillerNadathur2012} observe that such a language ``does not make
very strong use of the structure afforded by logic'' \book{75}.

This chapter lifts the restriction. The logic of \emph{first-order hereditary Harrop formulas},
\fohh, admits $D\imp G$ and $\forall x\,G$ in goals and hence, at any depth, in clause bodies. The
two new rules have a simple operational reading:
\begin{itemize}
  \item to prove $D\imp G$, prove $G$ with $D$ \emph{temporarily added to the program};
  \item to prove $\forall x\,G$, prove $G[c/x]$ with a \emph{new constant} $c$ temporarily added to
        the signature.
\end{itemize}
So ``programs and signatures can grow dynamically in the course of searching for a proof''
\book{75}, and they grow and shrink with a stack discipline. Three things follow, and they organise
the chapter: hypothetical reasoning becomes a programming technique (\cref{sec:hh-augment}); the
interplay of $\forall$ with logic variables imposes a \emph{scope condition} on unification
(\cref{sec:hh-generic,sec:hh-link}); and the declarative reading of the language is no longer
classical logic but intuitionistic, indeed minimal, logic (\cref{sec:hh-logic}). The full power of
$\forall$, binding \emph{names inside data}, ``must await the introduction of higher-order features''
\book{75}; it returns in \cref{ch:lambdatree}. Here $\forall$ shows its scoping face. The last section
(\cref{sec:hh-kanren}) collects what all this demands of a miniKanren, and of \micro{} in
particular.

%=====================================================================================================
\section{Syntax: goals, clauses and polarity}\label{sec:hh-syntax}

\begin{definition}[\fohh{} goals and clauses \book{75}]\label{def:hh-grammar}
Let $A$ range over first-order atomic formulas. Goals $G$ and program clauses $D$ are given by the
mutually recursive rules
\[
\begin{array}{rcl}
G &::=& \top \mid A \mid G\wedge G \mid G\vee G \mid \exists_\tau x\,G \mid D\imp G \mid \forall_\tau x\,G\\[2pt]
D &::=& A \mid G\imp D \mid D\wedge D \mid \forall_\tau x\,D
\end{array}
\]
A $D$-formula is called a \emph{first-order hereditary Harrop formula}; \fohh{} names both this class
of formulas and the logic programming framework built on it \book{76}.
\end{definition}

Compare with \fohc{} (\cref{ch:horn}), where $G::=\top\mid A\mid G\wedge G\mid G\vee G\mid\exists x\,G$
and $D$ has the same grammar. Two remarks make the definition easy to remember.
\begin{itemize}
  \item \emph{The $D$ grammar did not change.} The restrictions on clauses ``are, in fact, identical to
        the ones present in the \fohc{} setting'' \book{76}: no top-level $\vee$ or $\exists$, and the
        premise of a top-level implication is a goal. The class of clauses is larger only because goals,
        which appear as clause bodies, are now richer.
  \item \emph{The two grammars swap roles across $\imp$.} The premise of an implication \emph{goal} is a
        clause (it will be assumed); the premise of an implication \emph{clause} is a goal (it will be
        proved). Goals still are not freely generated: $(p\vee q)\imp r$ is not a goal, because
        $p\vee q$ is not a clause.
\end{itemize}

\paragraph{Concrete syntax.} In λProlog, $D\imp G$ in a goal is written \lp{D => G}; a clause
$G\imp A$ is written \lp{A :- G}; \lp{,} and \lp{&} both denote $\wedge$ (the chapter uses \lp{&}
inside assumed conjunctions of clauses), \lp{;} is $\vee$, \lp{pi x\ B} is $\forall x\,B$ and
\lp{sigma x\ B} is $\exists x\,B$. Free capitalised variables of a clause are implicitly universally
quantified at the front of the clause; free variables of a query are implicitly existential. The scope
of \lp{x\} extends as far to the right as possible.

\begin{pitfall}[the precedence of \texttt{=>}]
The connectives are listed in the order \lp{=>}, \lp{&}, \lp{,}, \lp{;}, \lp{:-}; their precedences
follow that order, so \lp{=>} binds \emph{tightest}; \lp{=>} and \lp{&} are right associative
\book{35}. Hence \lp{D => G1, G2} parses as \lp{(D => G1), G2}, and the hypothesis does \emph{not}
scope over \lp{G2}. This is why the parentheses in the queries of this chapter, such as
\lp{fact (finished X Y) => (fact (graduates X), fact (cs_major X))} and \lp{(p 2 & p 3) => p X}, are
essential, while \lp{p 2 => p 3 => p X} means $p\,2\imp(p\,3\imp p\,X)$.
\end{pitfall}

\subsection{Positive and negative occurrences}

With implications inside formulas it becomes meaningful to ask on which side of an arrow a
subformula sits.

\begin{definition}[polarity \book{76}]\label{def:hh-polarity}
$B$ is a positive subformula occurrence of $B$. If $C$ is a positive (negative) occurrence in $B$, then
it is a positive (negative) occurrence in $B\wedge B'$, $B'\wedge B$, $B\vee B'$, $B'\vee B$,
$B'\imp B$, $\forall x\,B$ and $\exists x\,B$, and it is a negative (positive) occurrence in $B\imp B'$.
Equivalently: $C$ is positive in $B$ if it lies to the left of an even number of implications, negative if
odd.
\end{definition}

\begin{proposition}[\book{76}]\label{prop:hh-polarity}
Positive subformulas of goals are goals and negative subformulas of goals are clauses; dually, positive
subformulas of clauses are clauses and negative ones are goals. A \fohc{} goal has no negative
subformulas. A hereditary Harrop formula is exactly a formula in which no positive subformula occurrence
is a disjunction or an existential quantification.
\end{proposition}

The last characterisation explains the name. The negative occurrences are the places where a formula
is \emph{used as a hypothesis}. Requiring that no positive part of a clause be a disjunction or an
existential says that whatever we assume, now or later in the search, never forces us to reason by cases
or to invent an anonymous witness. The adjective \emph{hereditary} refers to this
``at every positive position'' requirement (\cref{sec:hh-logic}).

\begin{example}
In the clause $\forall J\,((\forall x\,(\cst{bug}\,x\imp\cst{in}\,x\,J\imp\cst{dead}\,x))\imp
\cst{sterile}\,J)$ (\cref{sec:hh-generic}) the body $\forall x(\dots)$ is negative, hence a goal, and the
subformulas $\cst{bug}\,x$ and $\cst{in}\,x\,J$ are positive again (each lies to the left of exactly two
implications): they are clauses, and indeed they will be \emph{assumed}.
\end{example}

\subsection{Clausal order}

\begin{definition}[clausal order \book{76--77}]\label{def:hh-clausal}
\[
\begin{array}{r@{\;=\;}l}
\mathit{clausal}(A) & 0 \qquad (A\text{ atomic or }\top)\\
\mathit{clausal}(B_1\wedge B_2) = \mathit{clausal}(B_1\vee B_2) & \max(\mathit{clausal}(B_1),\mathit{clausal}(B_2))\\
\mathit{clausal}(B_1\imp B_2) & \max(\mathit{clausal}(B_1)+1,\ \mathit{clausal}(B_2))\\
\mathit{clausal}(\forall x\,B) = \mathit{clausal}(\exists x\,B) & \mathit{clausal}(B)
\end{array}
\]
\end{definition}

In \fohc{} goals have clausal order~0 and clauses order 0 or~1; in \fohh{} both are unbounded \book{77}.
Reading $\imp$ as the type arrow $\to$, the definition is the same as that of the order of a type
(\cref{ch:intro}); the book stresses that this is ``not accidental'' and will exploit it in its
Section~7.6 \book{77}. Worked values: \lp{q :- (q => q)}, that is $(q\imp q)\imp q$,
has order~2; the $\cst{sterile}$ clause above has order~2; the clause of \cref{fig:hh-rev3} below has
order~3.

%=====================================================================================================
\section{Operational semantics: AUGMENT and GENERIC}\label{sec:hh-rules}

Nothing new has to be invented. The framework of \cref{ch:horn} already contains the rules for
implications and universal quantifiers at the top of goals: ``These are to be treated by the AUGMENT and
GENERIC reduction rules'' \book{77}. Atomic goals are still solved by backchaining, which is possible
because $D$-formulas have the same shape as in \fohc{}, modulo the richer goals. For reference,
\cref{fig:hh-rules} collects the rules: the right-introduction rules for goals and the backchaining rules
for a clause $D$ ``in focus'', written $\hhfoc{\Sig}{\Prog}{D}{A}$.

\begin{figure}[t]
\small
\begin{mathpar}
\inferrule*[right=$\top$R]{ }{\seq{\Sig}{\Prog}{\top}}
\and
\inferrule*[right=$\wedge$R]{\seq{\Sig}{\Prog}{B_1} \\ \seq{\Sig}{\Prog}{B_2}}{\seq{\Sig}{\Prog}{B_1\wedge B_2}}
\and
\inferrule*[right=$\vee$R]{\seq{\Sig}{\Prog}{B_i}}{\seq{\Sig}{\Prog}{B_1\vee B_2}}
\and
\inferrule*[right=$\exists$R]{\seq{\Sig}{\Prog}{B[t/x]} \\ \Sig;\emptyset\vdash_f t:\tau}{\seq{\Sig}{\Prog}{\exists_\tau x\,B}}
\\
\inferrule*[right={$\imp$R\ (\rulename{augment})}]{\seq{\Sig}{\Prog,B_1}{B_2}}{\seq{\Sig}{\Prog}{B_1\imp B_2}}
\and
\inferrule*[right={$\forall$R\ (\rulename{generic})}]{\seq{c:\tau,\Sig}{\Prog}{B[c/x]}}{\seq{\Sig}{\Prog}{\forall_\tau x\,B}}\ \ c\notin\Sig
\\
\inferrule*[right=decide]{\hhfoc{\Sig}{\Prog}{D}{A}}{\seq{\Sig}{\Prog}{A}}\ \ D\in\Prog
\and
\inferrule*[right=initial]{ }{\hhfoc{\Sig}{\Prog}{A}{A}}
\and
\inferrule*[right=$\imp$L]{\hhfoc{\Sig}{\Prog}{D}{A} \\ \seq{\Sig}{\Prog}{G}}{\hhfoc{\Sig}{\Prog}{G\imp D}{A}}
\\
\inferrule*[right=$\wedge$L]{\hhfoc{\Sig}{\Prog}{D_i}{A}}{\hhfoc{\Sig}{\Prog}{D_1\wedge D_2}{A}}
\and
\inferrule*[right=$\forall$L]{\hhfoc{\Sig}{\Prog}{D[t/x]}{A} \\ \Sig;\emptyset\vdash_f t:\tau}{\hhfoc{\Sig}{\Prog}{\forall_\tau x\,D}{A}}
\end{mathpar}
\caption{The rules of O-proofs: right rules (book Figure 2.2) and backchaining rules (book Figure 2.3).
In \fohc{} the two rules of the second row are never used. With logic variables, $\exists$R and
$\forall$L introduce a new variable $X$ instead of guessing $t$, and initial unifies.}
\label{fig:hh-rules}
\end{figure}

\begin{definition}[O-proof \book{77}]\label{def:hh-oproof}
An \emph{O-proof} in \fohh{} is a derivation built with the typing rules for first-order terms and the
rules of \cref{fig:hh-rules} (book Figures 1.2, 2.2 and 2.3).
\end{definition}

In the logic-variable versions of the rules (\cref{ch:horn}, \book{61}) the typing premise of
$\exists$R and $\forall$L is delayed and initial unifies. \Cref{sec:hh-generic} shows that once
\textsc{generic} is present the delayed choice must still respect the signature in which the variable was
created.

\begin{tutorial}[programs grow and shrink like a stack]
Take the goal $(D_0\imp((D_1\imp G_1)\wedge(D_2\imp G_2)))\wedge G_3$ and the program $\Prog$ \book{77}.
\rulename{and} splits it in two. The left conjunct is an implication: \rulename{augment} moves $D_0$ into
the program, and \rulename{and} splits again; each conjunct is an implication again. In the end we must
prove
\[
G_1 \text{ from } \{D_1,D_0\}\cup\Prog,\qquad
G_2 \text{ from } \{D_2,D_0\}\cup\Prog,\qquad
G_3 \text{ from } \Prog.
\]
$D_1$ is never visible to $G_2$, and $D_0$ is never visible to $G_3$. Read as an interpreter state,
the program is pushed when we enter the antecedent's scope and popped when the subproof is done:
``implicational goals allow logic programs to grow and shrink following a stack-based discipline''
\book{78}. Notice that in the proof \emph{tree} nothing is popped: each branch simply carries its own
program. Popping is what a sequential interpreter does when it walks the tree.
\end{tutorial}

\paragraph{Programs are lists.} Under a depth-first strategy the order of clauses matters, so programs
are ``best thought of not as sets but rather as lists of clauses'' \book{78}. λProlog adds assumed
clauses \emph{at the front}: ``the most recently added clauses are the first to be used in backchaining''
\book{78}. Inside an assumed conjunction \lp{D1 & D2} the textual order is kept. With a program
containing only \lp{p 1}:
\begin{lprolog}
?- p 2 => p 3 => p X.
X = 3;
X = 2;
X = 1
?- (p 2 & p 3) => p X.
X = 2;
X = 3;
X = 1
\end{lprolog}

\begin{kanren}[answer order is a depth-first artefact]
Running the same two queries through the small interpreter of \cref{sec:hh-interp}, which uses the
interleaving \scm{or°} of \micro{} for the choice among clauses, gives the same \emph{sets} of answers
in a different order:
\begin{answers}
p 2 => p 3 => p X    : (3 2 1)
(p 2 & p 3) => p X   : (2 1 3)
\end{answers}
In the second query the $\wedge$L choice inside \lp{p 2 & p 3} is interleaved with the decide choice
over the rest of the program. Tests of this interpreter should not compare ordered answer lists with
λProlog's; the order of the program still matters, since front clauses are tried first. The proposed
module does reproduce the book's order here, \scm{(3 2 1)} and \scm{(2 3 1)}, because a dynamic relation
is one flat disjunction of its hypotheses, most recent first (\cref{def:design-a2},
\cref{sec:examples-order}). Compare ordered lists only where every clause answers without further search.
\end{kanren}

%=====================================================================================================
\section{Implicational goals and hypothetical reasoning}\label{sec:hh-augment}

\subsection{Entailment between Horn clauses}\label{sec:hh-entail}

Let $q,r,s,t,u$ be propositions (type $\typ{o}$) and consider the program \book{78}
\begin{lprolog}
s :- r, q.
t :- q, u.
q :- r.
\end{lprolog}
It has no atomic consequences at all, yet clauses follow from it: \lp{t :- r, u.} and \lp{s :- r.}
\book{79}. In \fohh{} we can ask this directly. The query \lp{?- s :- r.} ``is better written as''
\lp{?- r => s.} \book{79}: \rulename{augment} adds \lp{r} and ordinary Horn reasoning proves \lp{s}.

\begin{tutorial}[the O-proof of \texttt{r => s}]
Write $\Prog_0$ for the three clauses and $\Prog'=\Prog_0,r$; the signature is omitted. The search, read
goal-first (each line is reduced to the next ones):
\begin{enumerate}
  \item $\Prog_0\longrightarrow r\imp s$: by $\imp$R (\textsc{augment}) it suffices to show
        $\Prog'\longrightarrow s$.
  \item $\Prog'\longrightarrow s$: decide on the clause $r\wedge q\imp s$; by $\imp$L we need the head to
        match, $\Prog'\xrightarrow{\,s\,}s$ (initial), and the body, $\Prog'\longrightarrow r\wedge q$.
  \item $\Prog'\longrightarrow r\wedge q$: by $\wedge$R, both $\Prog'\longrightarrow r$ and
        $\Prog'\longrightarrow q$.
  \item $\Prog'\longrightarrow r$: decide on the assumed clause $r$, then initial.
  \item $\Prog'\longrightarrow q$: decide on $r\imp q$; initial for the head, and the body
        $\Prog'\longrightarrow r$ as in step~4.
\end{enumerate}
The only new step is the first one; the
rest is the Horn backchaining of \cref{ch:horn}, run in a larger program. In the same way
\lp{(r => u) => (r => t)} is provable \book{79}: it assumes the clauses \lp{r => u} and \lp{r}, and then
\lp{t} follows through \lp{q} and \lp{u}. (This is a different entailment from \lp{t :- r, u}, which as a
query is \lp{r => u => t}.)
\end{tutorial}

The last example of the section is a warning \book{79}:
\begin{lprolog}
?- (q :- (q => q)) => (q :- (q => q)).
\end{lprolog}
It has the form $B\imp B$, so it has an O-proof from the empty program, but a depth-first interpreter
``will not terminate on the given query''. To see why, unfold it. \rulename{augment} assumes
$(q\imp q)\imp q$; the goal is again $(q\imp q)\imp q$, so \rulename{augment} assumes $q\imp q$, i.e.\ the
clause \lp{q :- q}, \emph{at the front}. Now the program is \lp{[q :- q, q :- (q => q)]} and the goal
is \lp{q}. Depth-first search picks \lp{q :- q} and loops for ever. The O-proof picks the other clause:
its body \lp{q => q} assumes \lp{q} and proves \lp{q} at once.

\begin{pitfall}[provable, yet the interpreter loops]
Depth-first backchaining is incomplete, and AUGMENT makes this easier to hit, because assumed clauses
go to the front. The book makes the same point with conjunction: \lp{loop, fail} loops while
\lp{fail, loop} fails, with \lp{loop :- loop.} \book{88}. In the interleaving search of miniKanren the
clause choice is fair, and the interpreter of \cref{sec:hh-interp} does find the proof of $B\imp B$, in fact
infinitely many of them (\scm{run 3} returns \texttt{(α α α)}); conjunction, however, remains as
order-sensitive as in λProlog.
\end{pitfall}

\subsection{Hypothetical reasoning on a database}\label{sec:hh-db}

Implications in goals express \emph{hypothetical} questions: ``what if this were true?''.
\Cref{fig:hh-db} is the book's course database: facts about which students finished which courses, rules
for being a CS major and for graduating, and a constraint saying it is inconsistent to have finished both
210 and 250.

\begin{figure}[t]
\begin{lprolog}
kind entry               type.
type fact                entry -> o.
type false               o.
kind person              type.
type kim, dana           person.
type finished            person -> int -> entry.
type cs_major, graduates person -> entry.

fact (finished kim  102) & fact (finished dana 101).
fact (finished kim  210) & fact (finished dana 250).

fact (cs_major X)  :-
     (fact (finished X 101); fact (finished X 102)),
     fact (finished X 250), fact (finished X 301).
fact (graduates X) :-
     (fact (finished X 101); fact (finished X 102)),
     (fact (finished X 210); fact (finished X 250)),
     fact (finished X 301).

false :- fact (finished X 210), fact (finished X 250).
\end{lprolog}
\caption{An encoding of a small course-related database (book Figure 3.1, printed p.\,80).}
\label{fig:hh-db}
\end{figure}

Note the $D\wedge D$ clauses on the two fact lines, and that \lp{false} is an ordinary proposition: it
has a clause, so it can be proved, but it has no logical meaning (\cref{sec:hh-minimal}). The queries
and answers of the book \book{80--81}:
\begin{lprolog}
?- fact (finished dana X) => false.
X = 210;
no
?- fact (finished X Y) => (fact (graduates X), fact (cs_major X)).
X = dana
Y = 301;
no
?- fact (finished kim X) => fact (finished kim Y) =>
               (fact (graduates kim), fact (cs_major kim)).
X = 250
Y = 301
?- fact (finished kim 250) => fact (finished kim 301) => false.
solved
\end{lprolog}
The first asks whether some course would make the database inconsistent if Dana finished it (210, and
only 210). The second asks: ``If person X took just one more course Y, will he or she graduate with a CS
degree?'' (only Dana, with 301); this ``indirectly implies that there is no single course'' Kim can take
\book{80}, hence the third query, with two hypotheses: Kim needs 250 and 301. The fourth shows that this
plan is inconsistent, so ``there is no consistent way for Kim to graduate as a computer science major''
\book{81}.

\begin{tutorial}[a logic variable inside an assumed clause]
Follow the first query. Its free variable $X$ is existential, so $\exists$R makes it a logic variable
\emph{before} \rulename{augment} runs. The assumed clause is therefore \lp{fact (finished dana X)} with
$X$ \emph{not} universally quantified: it is one specific, still unknown fact. Now the goal \lp{false}
backchains on the constraint clause, whose own variable (call it $Z$) is renamed by $\forall$L:
\begin{enumerate}
  \item \lp{fact (finished Z 210)}: decide on the assumed clause unifies $Z=\cst{dana}$, $X=210$;
  \item \lp{fact (finished dana 250)}: true in the database.
\end{enumerate}
So proving \lp{false} \emph{binds the variable inside the hypothesis}, and the binding is reported to
the query. In the second query the single hypothesis \lp{fact (finished X Y)} likewise has one binding
for $X$ and $Y$, shared by every use of the clause. \Cref{sec:hh-link} turns this into a
programming technique.
\end{tutorial}

\begin{kanren}[the four queries on \micro]
The interpreter of \cref{sec:hh-interp} reproduces the four answers: \texttt{(210)} for Dana's
inconsistency, \texttt{((dana 301))} for the one-course question, a success for the inconsistency of Kim's
plan, and, asked for \emph{all} answers, \texttt{(250 301)} and the symmetric \texttt{(301 250)}, each
twice, for the two-course question: the book shows only the first answer, and the others correspond to
different proofs: the symmetric pair because either hypothesis can supply either course, the duplicates
because the ``210 or 250'' disjunct in the body of \lp{graduates} can be proved both from the database
(210) and from the hypothesis that supplies 250. As for what this demands of
the engine: the program must be part of what a goal sees, it must be extended for the duration of a
subgoal only, and an assumed clause must be able to \emph{mention logic variables of the surrounding
computation without renaming them} (\cref{sec:hh-kanren-context}).
\end{kanren}

\subsection{An interactive ``what if'' database}

\Cref{fig:hh-dbloop} turns hypothetical reasoning into a small interactive program; it uses the
non-logical \lp{print}, \lp{read} and cut \lp{!} \book{81}.

\begin{figure}[t]
\begin{lprolog}
type db                           o.
kind command                      type.
type do                           command -> o.
type enter, query, whatif, check  entry -> command.
type quit, consis                 command.

db :- print "Command?", read Command, do Command.

do quit.
do (enter Fact) :- fact Fact => db.
do (query Q) :- (fact Q, !, print "yes\n"; print "no\n"), db.
do (whatif Conjecture) :- (fact Conjecture => db),
                          print "Resuming\n", db.
do consis :- false, print "no\n", !; print "yes\n".
do (check Entry) :- (fact Entry, print "yes\n", !;
                    fact Entry => false, print "no\n", !;
                    print "no, but it could be true\n"),  db.
\end{lprolog}
\caption{A program that allows hypothetical queries against an encoded database (book Figure 3.2,
printed p.\,81).}
\label{fig:hh-dbloop}
\end{figure}

The book's session \book{81--82} is reproduced in \cref{fig:hh-session}.

\begin{figure}[t]
\noindent\begin{minipage}[t]{0.49\linewidth}
\begin{answers}
?- db.
Command? query (graduates dana).
no
Command? whatif (finished dana 301).
Command? query (graduates dana).
yes
Command? query (cs_major dana).
yes
Command? quit.
Resuming.
\end{answers}
\end{minipage}\hfill
\begin{minipage}[t]{0.49\linewidth}
\begin{answers}
Command? query (finished kim 101).
no
Command? query (graduates kim).
no
Command? whatif (finished kim 301).
Command? query (graduates kim).
yes
Command? query (cs_major kim).
no
Command? whatif (finished kim 250).
Command? query (cs_major kim).
yes
Command? consis.
no
Command?
\end{answers}
\end{minipage}
\caption{The book's session with the program of \cref{fig:hh-dbloop} (printed pp.\,81--82); read the
left column first.}
\label{fig:hh-session}
\end{figure}


The interesting point is where the ``database state'' lives: \emph{in the stack of nested
\rulename{augment}s}. \lp{enter F} runs the rest of the session inside \lp{fact F => db}, so $F$
persists; \lp{whatif F} runs a \emph{sub-session} inside \lp{fact F => db}, and when \lp{quit} ends that
sub-session the clause continues with \lp{print "Resuming\n", db} in the outer program, where $F$ is gone.
That is why Dana stops graduating after \lp{Resuming}. The Kim hypotheses nest: the innermost session sees
both 301 and 250, so \lp{cs_major kim} holds and \lp{consis} reports an inconsistency (210 and 250 together
prove \lp{false}). No \lp{assert}/\lp{retract} is involved: undoing a hypothesis is just leaving a scope.

\begin{remark}
The transcript is slightly idealised: \lp{do consis} is the only command besides \lp{quit} whose clause
does not end with \lp{db}, so after \lp{consis} the current sub-session ends and the enclosing \lp{whatif}
should print \texttt{Resuming} before the next prompt.
\end{remark}

%=====================================================================================================
\section{Universal goals and scoped constants}\label{sec:hh-generic}

\textsc{generic} transforms the goal $\forall_\tau x\,G(x)$ in the context $\langle\Sig,\Prog\rangle$ into
$G[c/x]$ in the context $\langle\Sig\cup\{c:\tau\},\Prog\rangle$ for a token $c$ not in $\Sig$; which token
is chosen does not matter \book{82}.

\begin{definition}[eigenvariable, scoped constant \book{82}]\label{def:hh-eigen}
The new token $c$ introduced by \textsc{generic} is an \emph{eigenvariable}. Since it is never
instantiated during computation, the book also calls it a \emph{scoped constant}.
\end{definition}

The parallel with implication is exact: implication augments the \emph{program} for part of the
computation, $\forall$ augments the \emph{signature} \book{82}.

\begin{tutorial}[sterile jars \book{83}]
A jar is sterile if every bug in it is dead; a bug in a heated jar is dead; jar \lp{j} has been heated:
\begin{lprolog}
kind jar, bug        type.
type j               jar.
type sterile, heated jar -> o.
type dead, bug       bug -> o.
type in              bug -> jar -> o.

sterile J :- pi x\ bug x => in x J => dead x.
dead B    :- heated J, in B J, bug B.
heated j.
\end{lprolog}
The signature has \emph{no constants of type $\typ{bug}$}: the program assumes nothing about the existence
of bugs. Now solve \lp{sterile j}.
\begin{enumerate}
  \item Backchain on the first clause: the goal becomes \lp{pi x\ bug x => in x j => dead x}.
  \item \textsc{generic}: pick a fresh \lp{g}; the goal becomes \lp{bug g => in g j => dead g}, and
        \lp{g : bug} joins the signature.
  \item \textsc{augment} twice: \lp{bug g} and \lp{in g j} join the program; the goal is \lp{dead g}.
  \item Backchain on the second clause: \lp{heated J} succeeds with $J=\cst{j}$, then \lp{in g j} and
        \lp{bug g} succeed by the two assumptions.
\end{enumerate}
``After this goal succeeds, the constant \lp{g} is removed from the signature, and the two clauses
\lp{bug g} and \lp{in g j} are similarly removed from the program'' \book{83}. The proof talks about
an arbitrary bug, never about a particular one.
\end{tutorial}

\subsection{An intensional universal quantifier}

The jar proof illustrates that \fohh's $\forall$ is \emph{intensional}: ``Proofs of universal goals do not
make any assumptions about the structure of the domain of quantification, and in fact, the same generic
proof must work for every element of the domain'' \book{83}. The alternative, \emph{extensional}, reading
would prove every instance, possibly in different ways, typically by induction; but ``the fohh logic does
not encompass rules for induction'' \book{83--84}. With
\begin{lprolog}
kind nat   type.
type zero  nat.
type succ  nat -> nat.
type plus  nat -> nat -> nat -> o.

plus zero L L.
plus (succ N) M (succ P) :- plus N M P.
\end{lprolog}
the query \lp{?- pi N\ plus N zero N.} is \emph{not} provable, although every closed instance is: the
generic goal \lp{plus c zero c} matches no clause head, since \lp{c} is neither \lp{zero} nor a
\lp{succ}. The base and step cases of the obvious induction are provable,
\begin{lprolog}
?- plus zero zero zero,
   pi N\ plus N zero N =>
         plus (succ N) zero (succ N).
\end{lprolog}
but no rule of \fohh{} combines them into \lp{pi N\ plus N zero N} \book{84}.

\begin{pitfall}[$\forall$ is not ``for every value'']
Do not expect \lp{pi} to enumerate a domain, and do not confuse it with the existential query
\lp{plus N zero N}, which enumerates \lp{zero}, \lp{succ zero}, \dots{} for ever. An empty type is harmless
for $\forall$: the jar proof works although no bug exists. The bibliographic notes call the two readings
\emph{generic} and \emph{universal} judgments; they coincide under the open-world assumption of \fohh,
and under a closed world Miller and Tiu's $\nabla$ quantifier captures the generic one \book{94},
\cite{MillerTiu2005}.

The dual remark concerns $\exists$. The typing premise $\Sig;\emptyset\vdash_f t:\tau$ of $\exists$R is,
with logic variables, \emph{delayed}: ``all the remaining, delayed typing judgments must be provable for
instances of the logic variables'' \book{61}. A query that ends with an unbound logic variable of an empty
type (say a \lp{sigma} over \lp{bug} in the jar signature, with a body that does not constrain it) has no
O-proof, because no closed term of type \lp{bug} exists, although an interpreter reports success. Inside
a \textsc{generic} scope over \lp{bug} the type is inhabited by the scoped constant, which is why the jar
proof is fine. An untyped miniKanren silently assumes every unbound answer variable has an instance.
\end{pitfall}

\subsection{Substitution must avoid capture}

In \fohc{} a clause is $\forall x_1\dots\forall x_m(A_1\wedge\dots\wedge A_n\imp A_0)$ with all quantifiers
outermost, so instantiating them is plain replacement \book{84}. In \fohh{} quantifiers occur inside
clauses, and plain replacement can capture.

\begin{definition}[free for \book{85}]\label{def:hh-freefor}
Let $x$ be a free variable of $B$ and $t$ a term of the same type. $t$ is \emph{free for} $x$ in $B$ if no
free occurrence of $x$ in $B$ is in the scope of a quantifier that binds a variable free in $t$.
\end{definition}
If $t$ is free for $x$ in $B$, replacement realises $B[t/x]$ soundly; otherwise one first renames bound
variables of $B$, and ``we shall assume that substitution is always performed in this logically sound
fashion'' \book{85}.

\begin{example}[\book{85}]
Substituting \lp{(f y)} for \lp{X} in \lp{p X :- pi y\ q X y} by naive replacement gives
\lp{p (f y) :- pi y\ q (f y) y}, where the constant \lp{y} has been captured: there are now two bound
occurrences where there was one. The correct result is \lp{p (f y) :- pi z\ q (f y) z}.
\end{example}

\subsection{Quantifier order matters}\label{sec:hh-scope}

``Possibly the simplest example illustrating the relevance of quantifier scopes in computation'' is
\book{85}
\begin{lprolog}
?- sigma x\ pi y\ x = y.
\end{lprolog}
It is not provable: \lp{pi y\ t = y} has no proof for any term $t$ in which $y$ does not occur. It is
\emph{true} in a model with a single element, but not every model is like that \book{85--86}. Proof
search agrees: \textsc{generic} creates a new constant $c$ and must prove $c=t$, which fails. Choosing
$y$ itself for $x$ does not help, because capture-avoiding substitution turns the goal into
\lp{pi z\ y = z} \book{86}. Swapping the quantifiers changes everything: \lp{pi y\ sigma x\ x = y}
\emph{is} provable, with witness the scoped constant itself.

\begin{tutorial}[the scope condition, derived from the rules]
Where does the asymmetry come from in the rules of \cref{fig:hh-rules}? From the side condition of
$\exists$R: the witness $t$ must satisfy $\Sig;\emptyset\vdash_f t:\tau$, where $\Sig$ is the signature
\emph{at the moment the existential is reduced}.
\begin{itemize}
  \item In $\exists x\,\forall y\,(x=y)$, $\exists$R is applied first, with $\Sig$ not containing the
        constant $c$ that \textsc{generic} will create later. No $\Sig$-term equals $c$.
  \item In $\forall y\,\exists x\,(x=y)$, \textsc{generic} runs first, the signature becomes $c,\Sig$, and
        $t=c$ is a legal witness.
\end{itemize}
With logic variables the witness is chosen \emph{later}, by unification. The side condition does not go
away; it becomes a constraint on unification: \emph{a logic variable created when the signature was
$\Sig$ may only be bound to terms built from $\Sig$} (and from other variables bound under the same
restriction). This is the scope condition behind the remark that unification ``must be modified \dots so
as to ensure that quantifier scopes are respected'' when \fohh{} computations alternate essential
existential and universal quantifiers \book{94}.
\end{tutorial}

The condition has a subtle consequence when a variable is bound to a term that contains \emph{other},
younger variables. Nadathur, Jayaraman and Kwon's implementation of $\forall$ goals tags every constant and
logic variable with a \emph{universe index}: \textsc{generic} increments the index and tags the new
constant with it, a logic variable is tagged with the index current at its creation, and $X$ (tag $i$) may
be bound to $t$ only if (1) the occurs check passes, (2) $t$ contains no constant with tag $>i$, and (3)
before binding, the tags of the variables of $t$ that are $>i$ are lowered to~$i$
\cite{NadathurJayaramanKwon1995}. Rule (3) is the subtle one. Their example: with the program clause
$\forall z\,(q(z)\imp p(d(z)))$, consider the goal $\exists x\,\forall y\,(q(y)\imp p(x))$.
\begin{enumerate}
  \item $\exists$: a logic variable $X$ with tag~1. $\forall$: a constant $c$ with tag~2. \textsc{augment}:
        assume $q(c)$; the goal is $p(X)$.
  \item Backchain on the clause: $\forall$L creates $Z$ with the current tag~2; unify $p(X)=p(d(Z))$,
        binding $X:=d(Z)$, which is allowed: $d$ belongs to the original signature (tag~1).
  \item The body $q(Z)$ must now be proved from the assumption $q(c)$, i.e.\ $Z:=c$. If $Z$ still had tag~2
        this would succeed, and $X$ would be $d(c)$, a term mentioning a constant that did not exist when
        $X$ was created. Rule (3) lowered $Z$ to tag~1 when $X$ was bound, so $Z:=c$ fails, as it should.
\end{enumerate}
The goal is not provable, since $x$ must be fixed before $y$, and only the lowering detects it.

\begin{kanren}[what $\forall$ demands of the unifier]
A miniKanren that wants \fohh's $\forall$ needs two things that \micro{} does not have
\cite{MaLuFriedman2020,HemannByrdFriedman2013,NadathurJayaramanKwon1995}:
\begin{enumerate}
  \item \emph{A new kind of term}, the scoped (eigen) constant: unique, never bound, equal only to
        itself. It must be distinct from every existing term, including records with equal slots,
        since \micro{} unifies records slot-wise.
  \item \emph{A scope check}: no variable that existed before the constant was created may ever reach it
        through the substitution, including through younger variables (rule (3) above). Known designs:
        an explicit quantifier prefix consulted by unification, as in λKanren 2020
        \cite{MaLuFriedman2020}; numeric universe levels with lowering, as in Teyjus and ELPI
        \cite{NadathurJayaramanKwon1995,DunchevEtAl2015}; an absento-like ``eigen'' constraint between the
        constants and the list of variables in scope, as in miniKanren with symbolic constraints
        \cite{HemannByrdFriedman2013}; or a check of the answers when the scope is left, as the
        interpreter of \cref{sec:hh-interp} does.
\end{enumerate}
Which of these fits \micro, and how variable creation numbers relate to levels, is a design question
taken up in \cref{ch:design}. Note also that αKanren's names do \emph{not} play this role: a nom ``only
unifies with itself or with an unassociated variable'' \cite{Byrd2009}, so nothing prevents an older
variable from being bound to it. \Cref{sec:hh-kanren-eigen} runs three tempting shortcuts and shows each
one proving $\exists x\,\forall y\,(x=y)$.
\end{kanren}

%=====================================================================================================
\section{Quantification links goals and clauses}\label{sec:hh-link}

In Prolog the program is fixed, so logic variables only ever occur in goals. In λProlog an
implication goal can assume a clause that mentions a logic variable created by an enclosing $\exists$ or by
the head of the clause being used; logic variables then ``may appear eventually in program clauses''
\book{86}. We met this in the Dana query. Here it becomes a technique.

\Cref{fig:hh-rev3} defines \lp{reverse} through an accumulator predicate \lp{rev} whose clauses exist only
while a \lp{reverse} goal is being solved.

\begin{figure}[t]
\begin{lprolog}
type reverse    list A -> list A -> o.
type rev        list A -> list A -> list A -> o.

reverse L K :-
   ((pi L\ rev nil L L) &
    (pi X\ pi L\ pi K\ pi M\ rev (X::L) K M :- rev L K (X::M)))
        => rev L K nil.
\end{lprolog}
\begin{lprolog}
type reverse, rev    list A -> list A -> o.

reverse L K :-
   (rev nil K &
    (pi X\ pi L\ pi K\ rev (X::L) K :- rev L (X::K)))
        => rev L nil.
\end{lprolog}
\caption{Two implementations of reverse with locally assumed clauses: with a three-place auxiliary
predicate (top, book Figure 3.3) and with a two-place one (bottom, book Figure 3.4), printed p.\,87.}
\label{fig:hh-rev3}
\label{fig:hh-rev4}
\end{figure}

With the first version, \lp{?- reverse (1::2::nil) P.} reduces to \lp{?- rev (1::2::nil) P nil.} in the
program augmented with
\begin{lprolog}
rev nil L L.
rev (X::L) K M :- rev L K (X::M).
\end{lprolog}
and eventually binds \lp{P} to \lp{(2::1::nil)} \book{86--87}. The explicit \lp{pi}s ``cannot be
dropped: Doing so will change the scope of the quantification and hence also the meaning'' \book{86}:
without \lp{pi L\} the clause \lp{rev nil L L} would mention the \emph{head's} \lp{L}, the list being
reversed, instead of being a general fact.

The second version starts from a different observation \book{87}. From the program
\begin{lprolog}
rv nil (c::b::a::nil).
rv (X::N) M :- rv N (X::M).
\end{lprolog}
one can prove \lp{rv (a::b::c::nil) nil}; replacing \lp{(c::b::a::nil)} by any list \lp{L}, the goal
\lp{rv K nil} is provable if and only if \lp{L} and \lp{K} are reverses of each other. This cannot be
written in \fohc, ``because to do so we require the ability to `tie' the binding for a variable in the
specification of the reverse predicate to one that appears in a goal at a particular point in the
computation'' \book{87}. In \fohh{} it is the second program of \cref{fig:hh-rev4}.

\begin{tutorial}[running the two-place reverse]
Solve \lp{?- reverse (1::2::nil) P.} with the second program.
\begin{enumerate}
  \item Backchain on the \lp{reverse} clause: $\forall$L renames its \lp{L}, \lp{K} to fresh logic
        variables; unification binds the first to \lp{1::2::nil} and identifies the second with \lp{P}.
  \item \textsc{augment} adds, at the front,
\begin{lprolog}
rev nil K.
rev (X::L) K :- rev L (X::K).
\end{lprolog}
        Here the \lp{K} of the first clause is the logic variable shared with \lp{P}: ``not like the
        usual implicitly universally quantified variables'' \book{88}. The \lp{K} of the second clause is
        bound by its own \lp{pi K\} and is renamed at every use.
  \item The goal \lp{rev (1::2::nil) nil} becomes \lp{rev (2::nil) (1::nil)}, then
        \lp{rev nil (2::1::nil)}.
  \item This last goal matches \lp{rev nil K} by binding the shared \lp{K}, ``and hence \lp{P}'', to
        \lp{2::1::nil} \book{88}.
\end{enumerate}
The base case of the recursion was written \emph{before} the recursion ran, with a hole for the answer;
the computation fills the hole from inside the scope.
\end{tutorial}

\begin{pitfall}[implicit quantification lies inside assumed clauses]
The display convention ``actually becomes somewhat misleading when such variables can appear in
clauses'' \book{88}: in \lp{rev nil K} the capital \lp{K} looks universally quantified but is a single
shared unknown. Two practical consequences. First, one may not copy an assumed clause wholesale at each
use (the \lp{copy_term} reflex of Prolog meta-interpreters): only the variables bound by the clause's own
quantifiers are renamed, the others stay shared \cite{NadathurJayaramanKwon1995}. Second, \lp{rev} is an
ordinary global predicate name, so the local clauses mix with any other assumption about \lp{rev}. Under a
stray assumption \lp{rev nil Z}, the interpreter of \cref{sec:hh-interp} returns two answers to the reverse
query, the second one produced by the stray clause:
\begin{answers}
under rev () Z       : (((2 1) α) (α (2 1)))
\end{answers}
Making \lp{rev} truly local requires quantifying over the predicate itself, a higher-order feature
(\cref{ch:hoprog}).
\end{pitfall}

\begin{kanren}[shared variables come for free with closures]
In a miniKanren embedded in Scheme, an assumed clause $\forall\bar x\,(G\imp p(\bar t))$ can be represented
by a host closure that, when called with the arguments of a goal, runs \scm{fresh°} for $\bar x$, unifies
the arguments with $\bar t$ and then runs $G$. Every use renames $\bar x$, as $\forall$L requires, while
every logic variable that is merely \emph{free} in the clause, like \lp{K} above or Dana's \lp{X}, is
captured by the closure and therefore shared, exactly the semantics of this section, with no copying
machinery. What closures do not provide is the scope information of \cref{sec:hh-scope} for the renamed
$\bar x$: in the universe-index scheme each such variable is tagged with the index current at the time of use
\cite{NadathurJayaramanKwon1995}, and that index must come from the search state. See \cref{sec:module-hh}.
\end{kanren}

%=====================================================================================================
\section{The logic behind fohh}\label{sec:hh-logic}

Logic is useful for programs because logically equivalent formulas, $B\equiv C$, ought to be
interchangeable. With side effects or incomplete search they are not: \lp{print "yes", print "no"}
differs from its swap, and \lp{loop, fail} from \lp{fail, loop}. ``Nevertheless \dots if we rule out
programs with side effects and consider complete proof-search strategies, then logical equivalences can be
used meaningfully in reasoning about logic programs'' \book{88}. The question is then: \emph{which} logic?

\subsection{Not classical logic}

For \fohc{} the rules are complete for both classical and intuitionistic logic (\cref{ch:horn}). Not so
for \fohh.

\begin{theorem}[\book{89}]\label{thm:hh-intuitionistic}
The operational semantics for sequents in \fohh{} is sound and complete for intuitionistic logic, but
classical logic is not sound for this operational semantics.
\end{theorem}

Take $p\vee(p\imp q)$ from the empty program. By $\vee$R it is provable iff $p$ is provable from $\emptyset$
(no) or $q$ is provable from $\{p\}$ (no). So there is no O-proof, in agreement with intuitionistic logic.
Classically, however \book{89},
\[
B_1\vee(B_2\imp B_3)\ \equiv\ B_1\vee\neg B_2\vee B_3\ \equiv\ (B_2\imp B_1)\vee B_3\ \equiv\ B_2\imp(B_1\vee B_3),
\]
so $p\vee(p\imp q)\equiv p\imp(p\vee q)$, which is clearly provable. Classical logic lets the scope of an
implication be ``extruded'' over a disjunction (the term comes from the π-calculus, \book{95}), and so
``the interpretation we desire for implication as a scoping mechanism is not valid in classical logic''
\book{89}. Three more classically valid goals without O-proofs \book{89}:
\[
(r\,a\wedge r\,b\imp q)\imp\exists x\,(r\,x\imp q)
\qquad\qquad
((p\imp q)\imp p)\imp p
\qquad\qquad
\exists x\,\forall y\,(p\,x\imp p\,y)
\]
The middle one is Peirce's formula. ``It is clear from all these examples that classical logic cannot
provide the declarative semantics of λProlog'' \book{89}.

\begin{tutorial}[why the three formulas fail]
\textbf{First.} \textsc{augment} assumes the clause \lp{q :- r a, r b}; $\exists$R makes $X$;
\textsc{augment} assumes $r\,X$; the goal $q$ backchains and needs $r\,a$ and $r\,b$. The only $r$-fact is
$r\,X$: the first subgoal binds $X:=a$ and then $r\,b$ fails. Classically one could argue ``by cases on
$x$''; a single O-proof must commit to one witness.

\textbf{Peirce.} \textsc{augment} assumes $(p\imp q)\imp p$, i.e.\ \lp{p :- (p => q)},
goal $p$. Backchaining gives the goal $p\imp q$: assume $p$, prove $q$. Nothing proves $q$.

\textbf{Drinker.} (The third formula is a form of the classical ``drinker paradox''; the name is not the
book's.) $\exists$R makes $X$, \textsc{generic} makes $c$, \textsc{augment} assumes $p\,X$, goal
$p\,c$. The only way is to unify $p\,X=p\,c$, i.e.\ $X:=c$, which the scope condition forbids. This is
exactly the failure of $\exists x\,\forall y\,(x=y)$, and indeed an engine without the scope check
``proves'' it (\cref{sec:hh-interp}).
\end{tutorial}

\subsection{Minimal and intuitionistic negation}\label{sec:hh-minimal}

Define $\neg B$ as $B\imp\bot$. \emph{Minimal logic} treats $\bot$ as essentially a non-logical constant;
\emph{intuitionistic logic} adds \emph{ex falso quodlibet} \book{89--90}:
\begin{mathpar}
\inferrule*[right=$\bot$R]{\seq{\Sig}{\Prog}{\bot}}{\seq{\Sig}{\Prog}{B}}
\end{mathpar}
A program is \emph{inconsistent} if it entails $\bot$; intuitionistically, it then entails everything
\book{90}. In λProlog one picks a proposition, say \lp{false}, for $\bot$; since it is non-logical, ex falso
is unavailable, and the negation obtained is \emph{minimal logic negation} \book{90}. That is what the
database constraint of \cref{fig:hh-db} uses. The interpreter can be extended to intuitionistic negation:
``before failing on a given query, the interpreter would try to determine if \lp{false} is provable. If it
is, then it would have to succeed on the query in question. Such a check for inconsistency would have to be
coordinated with each (sub)query that leads to the program being augmented'' \book{90}.

\begin{example}[where minimal and intuitionistic negation part ways]
Consider the goal \lp{p => (p => false) => q}, i.e.\ $p\imp\neg p\imp q$. \textsc{augment} twice puts
\lp{p} and the clause \lp{false :- p} in the program; the goal \lp{q} matches no clause head, so there is
no O-proof, and the formula is not provable in minimal logic either. If \lp{false} is read as $\bot$, it is
intuitionistically provable: \lp{false} follows by backchaining on \lp{false :- p}, and ex falso gives
\lp{q}. The extended interpreter described above finds this proof: before failing on \lp{q} it tries
\lp{false}, in the augmented program, and succeeds. The check must be made in the augmented program,
which is why it has to be ``coordinated'' with \textsc{augment}: the empty program is consistent, the
program $\{p,\ \neg p\}$ is not. (The interpreter of \cref{sec:hh-interp}, which has no ex falso, returns
\texttt{()} for this goal.)
\end{example}

Minimal negation is weak but not useless. These are provable \book{90}:
\begin{lprolog}
(p => q) => ((q => false) => (p => false)).
p => ((p => false) => false).
\end{lprolog}
Their converses are provable neither minimally nor intuitionistically, but are classically. The instance of
excluded middle \lp{?- p; (p => false).} is not provable; its double negation
\lp{?- ((p; (p => false)) => false) => false.} is.

\begin{tutorial}[a proof that uses a hypothesis twice \book{91}]
In the display below, the formulas assumed so far are written to the left of \lp{?-}:
\begin{lprolog}
(p; (p => false)) => false ?- false.
(p; (p => false)) => false ?- p; (p => false).
(p; (p => false)) => false ?- p => false.
p, (p; (p => false)) => false ?- false.
p, (p; (p => false)) => false ?- p; (p => false).
p, (p; (p => false)) => false ?- p.
\end{lprolog}
In line 1 \textsc{augment} has assumed the clause \lp{false :- p; (p => false)}. From line 1 to 2 we
backchain on it; from 2 to 3 we choose the right disjunct; from 3 to 4 \textsc{augment} assumes \lp{p};
from 4 to 5 we backchain on the \emph{same} clause again, now in a larger program; from 5 to 6 we choose the
left disjunct, and ``This last query succeeds immediately because the goal \lp{p} is also in the program at
that point.'' The pattern, a
hypothesis used twice, the second time in a bigger world, is typical of intuitionistic proofs.
\end{tutorial}

\begin{theorem}[\book{91}]\label{thm:hh-minimal}
If $\Prog$ is a \fohh{} program and $G$ a \fohh{} goal over $\Sig$, then $\seq{\Sig}{\Prog}{G}$ has an
O-proof if and only if it is provable in minimal logic.
\end{theorem}

Since \fohh{} formulas contain no $\bot$ (only the non-logical \lp{false}), minimal and intuitionistic
provability coincide on them, ``for this reason, we shall say that the declarative semantics of λProlog is
also given by intuitionistic logic'' \book{91}.

\begin{kanren}[\texttt{=> false} is not negation as failure]
Minimal negation needs nothing from the engine: \lp{false} is an ordinary predicate, and
$\neg D$ as a goal is ``assume $D$, prove \lp{false}''. It must not be confused with the soft-cut negation
\scm{(if° g ✗° ✓°)} available in \micro, which is negation as failure and unsound on non-ground goals, nor
with the disequality constraint \scm{≠°}. Under the open-world reading of \fohh{} the book notes that
$A\imp\bot$ for atomic $A$ ``can never be true because there is always some future world in which $A$ is
true'' \book{94}. Classical principles (excluded middle, Peirce, scope extrusion) must \emph{not} be built
into a \fohh{} engine: they are unsound for its operational semantics.
\end{kanren}

\subsection{Why ``hereditary Harrop''}

Harrop showed that if all assumptions (the formulas left of the sequent arrow) are of the form
$H::=\top\mid A\mid H\wedge H\mid B\imp H\mid\forall_\tau x\,H$, with $B$ arbitrary, then intuitionistic
logic has the \emph{disjunction property} (a proof of $F_1\vee F_2$ yields a proof of $F_1$ or of $F_2$)
and the \emph{existence property} (a proof of $\exists x\,F$ yields a witness $t$ with a proof of
$F[t/x]$) \cite{Harrop1960}, \book{93}. Intuitively, these are what goal-directed search relies on: when
the goal is a disjunction or an existential, committing to one disjunct or one witness loses nothing. Harrop's condition
gives them only at the root of a proof; to have them at every point of goal-directed search, disjunctions and existentials
are prohibited ``hereditarily at all positive locations'' \book{93}, whence hereditary Harrop formulas
\cite{MillerNadathurPfenningScedrov1991}. Semantically, intuitionistic logic is modelled by Kripke models
over possible worlds; identifying worlds with signature--program pairs, a \fohh{} computation moves to larger
worlds (the open-world assumption) \book{94}.

%=====================================================================================================
\section{Notable subsets of fohh}\label{sec:hh-subsets}

Can the grammar be simplified without losing expressiveness? The book gives two restrictions \book{91--92}:
\begin{align}
G &::= \top \mid A \mid G\wedge G \mid G\vee G \mid \exists x\,G \mid D\imp G \mid \forall x\,G
\qquad D ::= A \mid G\imp D \mid D\wedge D \mid \forall x\,D \tag{3.1}\\
D &::= A \mid G\imp A \mid D\wedge D \mid \forall x\,D \tag{3.2}\\
D &::= A \mid D\imp D \mid D\wedge D \mid \forall x\,D \tag{3.3}
\end{align}

\paragraph{Atomic heads, (3.2).} Clauses of (3.1) can be preprocessed, ``without changing their essential
operational interpretation'', into clauses whose top-level implications have atomic conclusions
\book{91}; declaratively, every (3.1) clause is intuitionistically (and minimally) equivalent to a
\emph{set} of (3.2) clauses \book{92}. The equivalences doing the work are those of \fohc{} normalisation
(\cref{ch:horn}, \book{64}):
\begin{gather*}
\forall x(B_1\wedge B_2)\equiv\forall x B_1\wedge\forall x B_2,\qquad
B_1\imp(B_2\imp B_3)\equiv(B_1\wedge B_2)\imp B_3,\\
B_1\imp(B_2\wedge B_3)\equiv(B_1\imp B_2)\wedge(B_1\imp B_3),
\end{gather*}
and $B_1\imp\forall x\,B_2\equiv\forall x(B_1\imp B_2)$ when $x$ is not free in $B_1$.

\begin{example}[normalising a clause to atomic heads]
The (3.1) clause $\forall x\,(g\,x\imp(p\,x\wedge(h\,x\imp q\,x)))$ has a conjunction and an implication
in the conclusion of its top-level implication. Distributing $\imp$ over $\wedge$, then $\forall$ over
$\wedge$, and finally currying the inner implication gives the two (3.2) clauses
\[
\forall x\,(g\,x\imp p\,x)\qquad\text{and}\qquad\forall x\,((g\,x\wedge h\,x)\imp q\,x),
\]
that is, \lp{p X :- g X.} and \lp{q X :- g X, h X.} Note that the body $g\,x$ now occurs twice. The book
warns, for \fohc, that this kind of preprocessing can make programs exponentially larger (from
$G\imp(D_1\wedge D_2)$ to $(G\imp D_1)\wedge(G\imp D_2)$ the goal $G$ is duplicated) and need not preserve
operational behaviour such as the number of derivations \book{64--65}. Which of $g\,x$ and $h\,x$ is
tried first in the second clause is an operational choice that the logic does not fix.
\end{example}

\paragraph{No $\vee$, no $\exists$, (3.3).} Dropping disjunction and existentials from goals gives
$G::=A\mid G\wedge G\mid D\imp G\mid\forall_\tau x\,G$ with $D$ as in (3.2); ``now the G- and D-formulas are
identical'' (up to the normalisation above), so one grammar (3.3) serves for both: atoms closed under
\emph{unrestricted} $\imp$, $\wedge$ and $\forall$ \book{92}.

\begin{pitfall}[equivalent as programs, not as formulas]
The passage from (3.1)/(3.2) to (3.3) is \emph{not} a logical equivalence \book{92}. The clause
$(p\imp(q\vee r))\imp s$ is legal in (3.2) but is not intuitionistically equivalent to any set of (3.3)
formulas; it implies $((p\imp q)\imp s)\wedge((p\imp r)\imp s)$, but not conversely (classically the two
are equivalent). Yet ``as program clauses, however, these two formulas prove exactly the same atomic
formulas'' \book{92}. The book does not spell out why; one way to see it is operational. To use the
clause on the goal $s$ we must prove the goal $p\imp(q\vee r)$, i.e.\ assume $p$ and then, in the
augmented program, commit to $q$ or to $r$: exactly what the two split clauses test. The intuitionistic
difference would only show if $p\imp(q\vee r)$ could be \emph{assumed}, and in \fohh{} it never is,
because it is not a clause.
\end{pitfall}

Another way to get rid of $\vee$ and $\exists$ is to define them as non-logical predicates with higher-order
Horn clauses (\cref{ch:hoprog}).

%=====================================================================================================
\section{What implication and universal goals ask of a miniKanren}\label{sec:hh-kanren}

We can now say precisely what is missing from \micro{} (\cref{ch:micro}). A \micro{} goal maps a state
(a substitution indexed by variable creation number, plus constraint stores) to a lazy stream of states;
relations are Scheme procedures, so \emph{the program is the fixed Scheme environment}: exactly the
\fohc{} situation this chapter started from.

\begin{center}
\small
\begin{tabularx}{\linewidth}{@{}lX@{}}
\toprule
\fohh & \micro{} today\\
\midrule
$\top$, $G\wedge G$, $G\vee G$ & \scm{✓°}, \scm{and°}, \scm{or°}/\scm{cond°} (fair interleaving)\\
$\exists x\,G$ ($\exists$R with a logic variable) & \scm{fresh°}\\
$\forall x\,D$ in a clause ($\forall$L) & \scm{fresh°} in a relation body\\
decide and $\wedge$L over the clauses of a predicate & the \scm{cond°} of a \scm{define-relation}\\
initial, by unification & \scm{=°} against the head\\
$D\imp G$ (\textsc{augment}) & missing: the program is not part of the state\\
$\forall x\,G$ (\textsc{generic}) & missing: no scoped constants, no scope check\\
\bottomrule
\end{tabularx}
\end{center}

\subsection{The program context belongs to the search state}\label{sec:hh-kanren-context}

The sequent $\seq{\Sig}{\Prog}{G}$ shows the program as an \emph{inherited, read-only} context: no rule
returns a modified program to its caller. In a miniKanren this translates into four requirements.
\begin{enumerate}
  \item \emph{A goal must be able to read the current program} when it reaches an atomic goal, so that
        assumed clauses are tried together with (and, in λProlog, before) the static ones.
  \item \emph{The extension is local to a subgoal.} After $G$ in $D\imp G$ produces an answer, the
        conjunction that follows must see the old program again. If the program lives in the state that
        \scm{and°} threads from one conjunct to the next, each answer state of $G$ must have its program
        restored.
  \item \emph{The extension is local to a branch.} Two disjuncts with different hypotheses must not see
        each other's clauses, even though the interleaving search alternates between them.
  \item \emph{Free logic variables of an assumed clause are shared}, not renamed
        (\cref{sec:hh-link}).
\end{enumerate}
Requirements 2 and 3 are what make a host-language \emph{dynamic binding} (\texttt{parameterize},
\texttt{fluid-let}, or a global variable set and restored) tempting, and wrong. A \micro{} goal does almost
nothing when applied to a state: \scm{fresh°} and \scm{define-relation} wrap the real work in a promise
(\scm{δ}), which is forced later by whoever consumes the stream, after the dynamic extent of the binding
has ended. The following file puts a ``current program'' in a parameter and tries it three ways.
\begin{scheme}[fontsize=\footnotesize]
(import scheme (chicken base) (only (scheme base) make-parameter parameterize)
        (aux base) (aux stream) (aux kanren micro))

;; The "current program" in a host-language parameter: a list of clauses (λ (x) goal).
(define program (make-parameter '()))

(define (any-clause° cs x)                     ; try every clause, most recent first
  (if (null? cs) ✗° (or° ((car cs) x) (any-clause° (cdr cs) x))))

(define-relation (hyp° x)                      ; a predicate whose clauses live in `program`
  (λ (s) ((any-clause° (program) x) s)))

(define (assume° clause g)                     ; D ⊃ G by dynamic binding: WRONG
  (λ (s) (parameterize ((program (cons clause (program)))) (g s))))

(define (fact a) (λ (y) (=° y a)))             ; the clause  hyp a.

(define (show x) (display x) (newline))
(show (μkanren-run (q -1 #t) (assume° (fact 'a) (hyp° q))))                  ; 1
(show (μkanren-run (q -1 #t) (parameterize ((program (list (fact 'a)))) (hyp° q)))) ; 2
(show (parameterize ((program (list (fact 'a)))) (μkanren-run (q -1 #t) (hyp° q)))) ; 3
\end{scheme}
\begin{answers}
()
()
(a)
\end{answers}
Run 1 binds the parameter around the \emph{application} of the goal to the state; run 2 around the
\emph{construction} of the goal. Both lose the hypothesis: \scm{(hyp° q)} returns a promise, and the
program is read only when the promise is forced, outside \texttt{parameterize}. Only run 3, which forces the
\emph{entire} search inside the binding, works, and it is useless for scoping, since the ``scope'' is then the
whole query. The parameter-based \texttt{assume} in the code accompanying λKanren 2020 loses hypotheses in
the same way for relation calls nested in delayed bodies; the paper also describes a state-threaded design
without the problem \cite{MaLuFriedman2020}. Hence: \emph{hypotheses must live in the search state} (or be
passed as an explicit argument, as in \cref{sec:hh-interp}), never in host dynamic state; see
\cref{ch:design}.

\subsection{Scoped constants need a scope check}\label{sec:hh-kanren-eigen}

Three shortcuts suggest themselves for $\forall y\,G$ in a miniKanren. Each one ``proves''
$\exists x\,\forall y\,(x=y)$, which has no proof (\cref{sec:hh-scope}):
\begin{scheme}[fontsize=\footnotesize]
(import scheme (chicken base) (aux base) (aux stream) (aux kanren micro))
(define (show x) (display x) (newline))

;; Three wrong encodings of the goal  ∃x ∀y (x = y),  which has no proof.

; (a) ∀ read as ∃: y is just another logic variable
(show (μkanren-run (x -1 #t) (fresh° (y) (=° x y))))

; (b) ∀ read as "a brand new constant", with no scope check
(show (μkanren-run (x -1 #t) (let ((c (gensym 'c))) (=° x c))))

; (c) ∀ read as "a fresh variable that must still be unbound at the end"
(define (naive-∀° f)
  (fresh° (y) (f y) (project° ((v y)) (if (μkanren-var? v) ✓° ✗°))))
(show (μkanren-run (x -1 #t) (naive-∀° (λ (y) (=° x y)))))
(show (μkanren-run (x -1 #t) (naive-∀° (λ (y) (=° x (list 'f y))))))
\end{scheme}
\begin{answers}
(α)
(c117)
(α)
((f α))
\end{answers}
(a) confuses the quantifiers. (b) gets a rigid constant but lets it escape into the answer: this is the
``scope extrusion'' that the scope condition forbids. (c) is the subtle one: \micro{} binds the older
variable \scm{x} to the younger \scm{y}, so \scm{y} stays unbound and the check passes; and with
\scm{(=° x (list 'f y))} it is \scm{x} that gets bound in \emph{any} miniKanren, so checking \scm{y} can
never detect the escape. A rigid constant \emph{plus} a check on the variables older than it is required,
as stated in the kanren box of \cref{sec:hh-scope}.

\begin{pitfall}[scoped constants and the constraint stores]
\micro{} also has disequality, type and absento constraints, and the book's language has none of them.
Moiseenko remarks that ``there is no sound implementation of eigen variables with the support of
disequality constraints'' \cite{Moiseenko2019}, and the \texttt{eigen} of miniKanren with symbolic
constraints is documented as unsupported with constraints other than \texttt{==}
\cite{HemannByrdFriedman2013}. Take $\forall x\,(x\neq\alpha)$, with $\alpha$ outer. Under an
extensional reading of $\forall$ \cite{JinEtAl2021} it should fail, since $x$ could be $\alpha$'s value;
under \rulename{generic}'s reading of $x$ as a new name, as in λProlog's \lp{pi x\ not (x = A)}, it holds.
A naive eigen constant gives the second answer by accident, without a scope check. The proposal implements
both readings consistently, the generic one by default (\cref{def:roadmap-readings},
\cref{sec:module-stores}).
\end{pitfall}

\subsection{A throwaway interpreter of the rules}\label{sec:hh-interp}

To check the examples of this chapter on today's \micro, and to see the rules of \cref{fig:hh-rules}
at work, here is a small interpreter of O-proofs. It is \emph{not} a proposal for the new module
(\cref{ch:design,ch:module} do that): formulas are data, it re-examines every clause at every atomic goal,
and the program is an explicit argument \scm{P}. But it is faithful to the rules, and it is complete with
respect to the clause choice thanks to the interleaving \scm{or°}.
\begin{scheme}[fontsize=\footnotesize]
(import scheme (chicken base) (aux base) (aux stream) (aux kanren micro))

;; A throwaway interpreter of O-proofs over (aux kanren micro): formulas are data,
;; binders are Scheme closures.
;;   G ::= #t | A | (and G ...) | (or G ...) | (some ,(λ (x) G)) | (eq t u)
;;       | (imp D G) | (all ,(λ (x) G))
;;   D ::= A | (if G D) | (both D D) | (all ,(λ (x) D))
(define (tag? t F) (and (pair? F) (eq? (car F) t)))

(define (all-of° gs) (if (null? gs) ✓° (and° (car gs) (μ s ((all-of° (cdr gs)) s)))))
(define (one-of° gs) (if (null? gs) ✗° (or° (car gs) (μ s ((one-of° (cdr gs)) s)))))

(define (solve° P G)                                  ; Σ; P ⟶ G
  (cond
    ((eq? G #t)       ✓°)                                           ; ⊤R
    ((tag? 'and G)    (all-of° (map (λ (g) (solve° P g)) (cdr G))))  ; ∧R
    ((tag? 'or G)     (one-of° (map (λ (g) (solve° P g)) (cdr G))))  ; ∨R
    ((tag? 'some G)   (fresh° (x) (solve° P ((cadr G) x))))          ; ∃R, logic variable
    ((tag? 'eq G)     (=° (cadr G) (caddr G)))
    ((tag? 'imp G)    (solve° (cons (cadr G) P) (caddr G)))         ; AUGMENT, at the front
    ((tag? 'all G)    (generic° P (cadr G) #t))                     ; GENERIC
    ((tag? 'all/unchecked G) (generic° P (cadr G) #f))
    (else             (decide° P P G))))                            ; atomic goal

(define (decide° P Ds A)                              ; decide: some D ∈ P, front first
  (if (null? Ds) ✗° (or° (bc° P (car Ds) A) (μ s ((decide° P (cdr Ds) A) s)))))

(define (bc° P D A)                                   ; Σ; P —D→ A
  (cond
    ((tag? 'if D)   (and° (bc° P (caddr D) A) (μ s ((solve° P (cadr D)) s))))  ; ⊃L
    ((tag? 'both D) (or° (bc° P (cadr D) A) (bc° P (caddr D) A)))            ; ∧L
    ((tag? 'all D)  (fresh° (x) (bc° P ((cadr D) x) A)))                      ; ∀L
    (else           (=° D A))))                                               ; initial

;; GENERIC: a new constant c; an answer of G[c/x] survives only if no variable
;; created before c (index < n) reaches c through the substitution.
(define (mentions? c w)                              ; does the walked term w contain c?
  (or (eq? w c) (and (pair? w) (or (mentions? c (car w)) (mentions? c (cdr w))))))

(define (generic° P body check?)
  (λ (s)
    (let* ((n (μkanren-state-vars-count s))
           (c (gensym 'c))
           (ok? (λ (s*) (let loop ((i 0))
                          (or (= i n)
                              (and (not (mentions? c (μkanren-state-find/value (make-μkanren-var i) s*)))
                                   (loop (add1 i)))))))
           (answers ((solve° P (body c)) s)))
      (if check? (filter§ ok? answers) answers))))

(define-syntax ask                                    ; (ask n (q) P G): first n answers, -1 = all
  (syntax-rules () ((_ n (q) P G) (μkanren-run (q n #t) (solve° P G)))))
(define (show label x) (display label) (display x) (newline))
\end{scheme}

How the code mirrors the rules: \textsc{augment} is \scm{(solve° (cons D P) G)}, so the new clause goes to
the front, the scope ends with the Scheme call, and each branch carries its own \scm{P} (requirements 2
and 3 of \cref{sec:hh-kanren-context} hold by construction). Binders are Scheme closures, so $\forall$L and
$\exists$R call \scm{fresh°} when the clause or goal is used: a clause's own variables are renamed at each
use, while a logic variable merely mentioned in it is shared (requirement 4). \textsc{generic} creates a
\scm{gensym} constant (\micro{} compares symbols with \scm{eq?}) and enforces the scope condition lazily:
an answer of $G[c/x]$ is kept only if no variable with creation index below \scm{n}, i.e.\ older than
\scm{c}, reaches \scm{c}; the deep walk also catches escapes through younger variables. It prunes late,
which is acceptable for a test bench, not for an engine. \scm{all/unchecked} omits the check.

\paragraph{The examples.} The file below begins with \scm{(include "hh.scm")}, the listing above, and
runs a selection of the chapter's queries, in the book's order.
\begin{scheme}[fontsize=\footnotesize]
(include "hh.scm")
;; Section 3.2 of the book: clause order, Horn entailment, B ⊃ B
(show "p 2 => p 3 => p X    : " (ask -1 (x) '((p 1)) `(imp (p 2) (imp (p 3) (p ,x)))))
(show "(p 2 & p 3) => p X   : " (ask -1 (x) '((p 1)) `(imp (both (p 2) (p 3)) (p ,x))))
(define P0 '((if (and r q) s) (if (and q u) t) (if r q)))
(show "r => s ; t           : " (list (ask -1 (q) P0 '(imp r s)) (ask -1 (q) P0 't)))
(show "B => B, 3 proofs     : " (ask 3 (q) '() '(imp (if (imp q q) q) (imp (if q q) q))))
;; Figure 3.1, the course database
(define (fin x c) `(fact (finished ,x ,c)))
(define DB
  `((both ,(fin 'kim 102) ,(fin 'dana 101))
    (both ,(fin 'kim 210) ,(fin 'dana 250))
    (all ,(λ (X) `(if (and (or ,(fin X 101) ,(fin X 102)) ,(fin X 250) ,(fin X 301))
                      (fact (cs_major ,X)))))
    (all ,(λ (X) `(if (and (or ,(fin X 101) ,(fin X 102)) (or ,(fin X 210) ,(fin X 250))
                           ,(fin X 301))
                      (fact (graduates ,X)))))
    (all ,(λ (X) `(if (and ,(fin X 210) ,(fin X 250)) false)))))
(show "dana X => false      : " (ask -1 (x) DB `(imp ,(fin 'dana x) false)))
(show "X Y => grad, major   : "
      (μkanren-run (q -1 #t)
        (fresh° (x y) (=° q (list x y))
          (solve° DB `(imp ,(fin x y) (and (fact (graduates ,x)) (fact (cs_major ,x))))))))
(show "kim X => kim Y => .. : "
      (μkanren-run (q -1 #t)
        (fresh° (x y) (=° q (list x y))
          (solve° DB `(imp ,(fin 'kim x) (imp ,(fin 'kim y)
                         (and (fact (graduates kim)) (fact (cs_major kim)))))))))
(show "kim 250, 301 => false: " (ask -1 (q) DB `(imp ,(fin 'kim 250) (imp ,(fin 'kim 301) false))))
;; Section 3.3: jars, quantifier order, levels
(define JAR
  `((all ,(λ (J) `(if (all ,(λ (x) `(imp (bug ,x) (imp (in ,x ,J) (dead ,x))))) (sterile ,J))))
    (all ,(λ (B) `(all ,(λ (J) `(if (and (heated ,J) (in ,B ,J) (bug ,B)) (dead ,B))))))
    (heated j)))
(show "sterile Q            : " (ask -1 (q) JAR `(sterile ,q)))
(show "pi y\\ X = f y        : " (ask -1 (x) '() `(all ,(λ (y) `(eq ,x (f ,y))))))
(show "  ... unchecked      : " (ask -1 (x) '() `(all/unchecked ,(λ (y) `(eq ,x (f ,y))))))
(show "pi y\\ sigma x\\ x = y : " (ask -1 (q) '() `(all ,(λ (y) `(some ,(λ (x) `(eq ,x ,y)))))))
(define PQ `((all ,(λ (z) `(if (q ,z) (p (d ,z)))))))
(show "pi y\\ (q y => p X)   : " (ask -1 (x) PQ `(all ,(λ (y) `(imp (q ,y) (p ,x))))))
(show "  ... unchecked      : " (ask -1 (x) PQ `(all/unchecked ,(λ (y) `(imp (q ,y) (p ,x))))))
;; Figure 3.4: reverse; the base clause captures the outer K
(define REV
  `((all ,(λ (L) `(all ,(λ (K)
      `(if (imp (both (rev () ,K)
                      (all ,(λ (X) `(all ,(λ (L) `(all ,(λ (K)
                          `(if (rev ,L (,X . ,K)) (rev (,X . ,L) ,K)))))))))
                (rev ,L ()))
           (reverse ,L ,K))))))))
(show "reverse (1 2) P      : " (ask -1 (p) REV `(reverse (1 2) ,p)))
(show "reverse P (1 2 3)    : " (ask 1 (p) REV `(reverse ,p (1 2 3))))
(show "under rev () Z       : "
      (μkanren-run (q -1 #t)
        (fresh° (p z) (=° q (list p z)) (solve° REV `(imp (rev () ,z) (reverse (1 2) ,p))))))
;; Section 3.4: Peirce, the drinker, double negation of excluded middle
(show "Peirce               : " (ask -1 (q) '() '(imp (if (imp p q) p) p)))
(show "drinker              : " (ask -1 (q) '() `(some ,(λ (x) `(all ,(λ (y) `(imp (p ,x) (p ,y))))))))
(show "  ... unchecked      : "
      (ask -1 (q) '() `(some ,(λ (x) `(all/unchecked ,(λ (y) `(imp (p ,x) (p ,y))))))))
(show "p ; (p => false)     : " (ask -1 (q) '() '(or p (imp p false))))
(show "not not (p ; not p)  : " (ask 1 (q) '() '(imp (if (or p (imp p false)) false) false)))
(show "p => not p => q      : " (ask -1 (q) '() '(imp p (imp (if p false) q))))
\end{scheme}
\begin{answers}
p 2 => p 3 => p X    : (3 2 1)
(p 2 & p 3) => p X   : (2 1 3)
r => s ; t           : ((α) ())
B => B, 3 proofs     : (α α α)
dana X => false      : (210)
X Y => grad, major   : ((dana 301))
kim X => kim Y => .. : ((301 250) (250 301) (250 301) (301 250))
kim 250, 301 => false: (α)
sterile Q            : (j)
pi y\ X = f y        : ()
  ... unchecked      : ((f c488))
pi y\ sigma x\ x = y : (α)
pi y\ (q y => p X)   : ()
  ... unchecked      : ((d c563))
reverse (1 2) P      : ((2 1))
reverse P (1 2 3)    : ((3 2 1))
under rev () Z       : (((2 1) α) (α (2 1)))
Peirce               : ()
drinker              : ()
  ... unchecked      : (α)
p ; (p => false)     : ()
not not (p ; not p)  : (α)
p => not p => q      : ()
\end{answers}
A provable closed goal answers \texttt{(α)} (a proof, nothing to report); an unprovable one answers
\texttt{()}. Every line agrees with the book; the Kim query lists all answers, while the book shows only
the first. The goal \texttt{t} on its own fails, since the Horn program has no atomic consequences, and the
last line is the minimal-logic failure of $p\imp\neg p\imp q$ discussed in \cref{sec:hh-minimal}. In \texttt{sterile Q} the outer variable is bound to the constant \lp{j}, which is fine: \lp{j}
belongs to the signature in which \lp{Q} was created. In the reverse example, \scm{(rev () ,K)} is built
with the \emph{outer} \scm{K}, a logic variable, which the closure captures: the base case shares the
answer with the query. The three \texttt{unchecked} lines are the instructive ones. Without the scope
condition a scoped constant leaks into an answer, once directly and once through the clause variable $Z$ of
the level-lowering example, and the drinker formula $\exists x\,\forall y\,(p\,x\imp p\,y)$, valid
classically but not intuitionistically, becomes ``provable''. Dropping the scope check does not only produce
odd answers: it changes the logic.

%=====================================================================================================
\section{Summary}\label{sec:hh-summary}

\begin{itemize}
  \item \fohh{} extends \fohc{} with $D\imp G$ and $\forall x\,G$ in goals; the clause grammar is unchanged,
        and the two grammars swap roles across $\imp$. Hereditary Harrop formulas are those with no positive
        $\vee$ or $\exists$ (\cref{def:hh-grammar,prop:hh-polarity}).
  \item \textsc{augment} proves $D\imp G$ with $D$ added at the front of the program; \textsc{generic} proves
        $\forall x\,G$ with a new scoped constant in the signature. Both follow a stack discipline
        (\cref{fig:hh-rules}).
  \item Implication gives hypothetical reasoning: entailment between clauses, ``what if'' queries, and
        a database whose state is the stack of open hypotheses (\cref{sec:hh-augment}).
  \item $\forall$ is intensional: one generic proof for all elements, no induction, empty types harmless
        (\cref{sec:hh-generic}).
  \item With logic variables, $\exists\forall$ and $\forall\exists$ differ through a \emph{scope condition}: a
        variable may only be bound to terms over the signature of its creation, and binding must propagate this
        restriction to the variables of the bound term (universe levels and lowering, \cref{sec:hh-scope}).
  \item Assumed clauses may contain logic variables of the surrounding computation; these are shared, not
        renamed, which lets a clause ``tie'' its binding to a goal (reverse, \cref{sec:hh-link}).
  \item O-proofs are sound and complete for intuitionistic logic and coincide with minimal provability on
        \fohh; classical logic is unsound for them (scope extrusion, Peirce, drinker). \lp{false} gives
        minimal negation, not negation as failure (\cref{thm:hh-intuitionistic,thm:hh-minimal}).
  \item Clauses can be normalised to atomic heads (3.2); removing $\vee$ and $\exists$ (3.3) preserves atomic
        consequences but not logical equivalence (\cref{sec:hh-subsets}).
  \item For a miniKanren: the program must live in the search state (per branch, restored after each answer of
        the scoped subgoal), never in host dynamic binding, because goals run lazily; $\forall$ needs rigid
        scoped constants and a scope check in, or around, unification (\cref{sec:hh-kanren}, and
        \cref{ch:design} for the design).
\end{itemize}
